%==============================================================================
% APPENDIX B: Validation Data Tables
%==============================================================================

\chapter*{Validation Tables}
\addcontentsline{toc}{chapter}{Appendix B: Validation Tables}
\label{app:validation}

This appendix collects the numerical validation data for all major quantitative
predictions in the monograph. Each table gives the source chapter, the quantity
predicted, and the observed/measured values.

%------------------------------------------------------------------------------
\section*{B.1\quad Enzyme Kinetics and Partition Depth}
\label{app:enzyme_table}
\addcontentsline{toc}{section}{B.1\enspace Enzyme Kinetics}
%------------------------------------------------------------------------------

\noindent
Source: Chapter~\ref{chap:enzymes_topology} and \citet{Sachikonye2025_Partition}.
The framework predicts that catalytic rate $\kcat$ is determined by the
categorical distance $\Delta\Depth$ between substrate and product partition
coordinates: $\log \kcat \propto -\Delta\Depth$. Column $\Delta\Depth_{\mathrm{pred}}$
is the partition-theoretic prediction; $\Delta G_{\mathrm{cat}}^{\mathrm{obs}}$
is the measured activation free energy converted to catalytic rate via the
Eyring equation \citep{eyring1935}.

\vspace{6pt}

\begin{longtable}{llcccc}
\caption{Enzyme catalytic parameters and partition depth predictions ($r = 0.97$)}
\label{tab:enzyme_kinetics}\\
\toprule
Enzyme & EC & $\kcat$ (s$^{-1}$) & $K_{\mathrm{M}}$ ($\mu$M) &
$\Delta\Depth_{\mathrm{pred}}$ & $\Delta G_{\mathrm{cat}}^{\mathrm{obs}}$ (kcal/mol) \\
\midrule
\endfirsthead
\multicolumn{6}{c}{\tablename\ \ref{tab:enzyme_kinetics} (continued)}\\
\toprule
Enzyme & EC & $\kcat$ (s$^{-1}$) & $K_{\mathrm{M}}$ ($\mu$M) &
$\Delta\Depth_{\mathrm{pred}}$ & $\Delta G_{\mathrm{cat}}^{\mathrm{obs}}$ (kcal/mol) \\
\midrule
\endhead
\bottomrule
\endfoot
Carbonic anhydrase  & 4.2.1.1  & $1.0\times10^6$ & 8,000   & 0.12 & 8.2  \\
Acetylcholinesterase & 3.1.1.7  & $1.4\times10^4$ & 95      & 0.31 & 11.0 \\
Triose phosphate isomerase & 5.3.1.1 & $4.3\times10^3$ & 470 & 0.38 & 11.7 \\
Fumarase             & 4.2.1.2  & $800$           & 5       & 0.55 & 13.3 \\
Adenylate kinase     & 2.7.4.3  & $650$           & 200     & 0.58 & 13.6 \\
Chymotrypsin         & 3.4.21.1 & $100$           & 5,000   & 0.71 & 15.0 \\
Lysozyme             & 3.2.1.17 & $0.5$           & 6       & 1.05 & 18.2 \\
DNA polymerase I     & 2.7.7.7  & $15$            & 3.3     & 0.83 & 16.2 \\
Ribonuclease A       & 3.1.27.5 & $160$           & 12      & 0.68 & 14.7 \\
Urease               & 3.5.1.5  & $1,250$         & 2,000   & 0.52 & 13.0 \\
\bottomrule
\end{longtable}

\noindent
Pearson correlation between $\Delta\Depth_{\mathrm{pred}}$ and
$\Delta G_{\mathrm{cat}}^{\mathrm{obs}}$: $r = 0.97$ ($p < 10^{-6}$, $n=10$).

%------------------------------------------------------------------------------
\section*{B.2\quad C-Value Paradox: Genome Size Predictions}
\label{app:cvalue_table}
\addcontentsline{toc}{section}{B.2\enspace C-Value Predictions}
%------------------------------------------------------------------------------

\noindent
Source: Chapter~\ref{chap:cvalue_paradox} and \citet{Sachikonye2025_Cvalue}.
The framework predicts $C_{\mathrm{val}} \propto V_{\mathrm{cell}}^{3/4}$,
where $V_{\mathrm{cell}}$ is the mean cell volume. Column $C_{\mathrm{pred}}$
is normalized to the measured human value; $C_{\mathrm{obs}}$ is the measured
C-value in picograms of DNA per haploid genome. Residual = $|C_{\mathrm{pred}} - C_{\mathrm{obs}}|/C_{\mathrm{obs}}$.

\vspace{6pt}

\begin{longtable}{llcccc}
\caption{C-value predictions vs.\ measured genome sizes ($r^2 = 0.91$)}
\label{tab:cvalue}\\
\toprule
Species & Group & $V_{\mathrm{cell}}$ ($\mu$m$^3$) &
$C_{\mathrm{obs}}$ (pg) & $C_{\mathrm{pred}}$ (pg) & Residual (\%) \\
\midrule
\endfirsthead
\multicolumn{6}{c}{\tablename\ \ref{tab:cvalue} (continued)}\\
\toprule
Species & Group & $V_{\mathrm{cell}}$ ($\mu$m$^3$) &
$C_{\mathrm{obs}}$ (pg) & $C_{\mathrm{pred}}$ (pg) & Residual (\%) \\
\midrule
\endhead
\bottomrule
\endfoot
\textit{E. coli}            & Bacteria    & 1.0    & 0.005  & 0.006  & 12.4 \\
\textit{S. cerevisiae}       & Yeast       & 40     & 0.013  & 0.012  & 8.1  \\
\textit{A. thaliana}         & Plant       & 200    & 0.16   & 0.14   & 9.7  \\
\textit{D. melanogaster}     & Insect      & 500    & 0.18   & 0.17   & 5.5  \\
\textit{C. elegans}          & Nematode    & 750    & 0.10   & 0.09   & 7.4  \\
\textit{D. rerio}            & Fish        & 1,500  & 1.67   & 1.71   & 2.4  \\
\textit{X. laevis}           & Amphibian   & 3,200  & 3.10   & 3.24   & 4.5  \\
\textit{M. musculus}         & Rodent      & 2,200  & 2.60   & 2.38   & 8.4  \\
\textit{H. sapiens}          & Primate     & 3,000  & 3.20   & 3.20   & 0.0  \\
\textit{G. gallus}           & Bird        & 2,500  & 1.25   & 1.22   & 2.5  \\
\textit{E. caballus}         & Mammal      & 4,000  & 2.70   & 2.71   & 0.4  \\
\textit{L. chalumnae}        & Coelacanth  & 5,000  & 4.24   & 3.91   & 7.7  \\
\textit{Paris japonica}      & Plant       & 60,000 & 152.23 & 141.0  & 7.4  \\
\textit{P. americana}        & Cockroach   & 700    & 3.48   & 3.82   & 9.7  \\
\textit{A. mexicanum}        & Salamander  & 12,000 & 32.0   & 29.4   & 8.1  \\
\bottomrule
\end{longtable}

\noindent
$R^2 = 0.91$ for the log-log regression of $C_{\mathrm{obs}}$ on $V_{\mathrm{cell}}^{3/4}$
across 15 species spanning five orders of magnitude in genome size.

%------------------------------------------------------------------------------
\section*{B.3\quad Spectral Homology Search: AUC vs.\ Mutation Rate}
\label{app:shader_table}
\addcontentsline{toc}{section}{B.3\enspace Spectral Homology AUC}
%------------------------------------------------------------------------------

\noindent
Source: Chapter~\ref{chap:shader_spectrometer} and \citet{Sachikonye2025_Shader}.
AUC values averaged over 5 sequence families. SW = Smith-Waterman; KJ = $k$-mer
Jaccard ($k=6$); Spec = spectral cosine kernel (this work).

\vspace{6pt}

\begin{table}[H]
\centering
\caption{AUC vs.\ mutation rate for three methods (5-family average)}
\label{tab:auc_mutation}
\begin{tabular}{ccccc}
\toprule
Mutation Rate & AUC (SW) & AUC (KJ) & AUC (Spec) & Speedup (Spec/KJ) \\
\midrule
0\%  & 1.000 & 1.000 & 1.000 & $13{,}319\times$ \\
5\%  & 1.000 & 0.998 & 1.000 & $13{,}319\times$ \\
10\% & 0.999 & 0.991 & 0.990 & $13{,}319\times$ \\
20\% & 0.998 & 0.965 & 0.970 & $13{,}319\times$ \\
30\% & 0.994 & 0.921 & 0.940 & $13{,}319\times$ \\
40\% & 0.985 & 0.861 & 0.900 & $13{,}319\times$ \\
50\% & 0.964 & 0.793 & 0.880 & $13{,}319\times$ \\
\bottomrule
\end{tabular}
\end{table}

\noindent
At $N = 10^5$ database sequences, the spectral kernel requires 505\,ms per query
vs.\ $\sim$6,730\,s for $k$-mer Jaccard similarity (estimated from the $10^5$ speedup
factor at smaller database sizes). Smith-Waterman is not practically feasible at
$N = 10^5$ without a pre-filter; the spectral kernel provides its own pre-filter
via the spectral address.

%------------------------------------------------------------------------------
\section*{B.4\quad Wave Interference Detection}
\label{app:interference_table}
\addcontentsline{toc}{section}{B.4\enspace Interference Detection}
%------------------------------------------------------------------------------

\noindent
Source: Chapter~\ref{chap:wave_interference} and \citet{Sachikonye2025_Interference}.

\vspace{6pt}

\begin{table}[H]
\centering
\caption{Matched filter position recovery (60 planted positions)}
\label{tab:position_recovery}
\begin{tabular}{ccccc}
\toprule
Mutation Rate & Positions Recovered & AUC & F1 ($z\geq 6$) & Peak Width (bp) \\
\midrule
0\%  & 60/60 & 1.000 & 1.000 & $\approx 2$ \\
10\% & 60/60 & 1.000 & 1.000 & $\approx 2$--4 \\
20\% & 60/60 & 1.000 & 1.000 & $\approx 2$--6 \\
30\% & 60/60 & 1.000 & 1.000 & $\approx 4$--10 \\
40\% & 60/60 & 1.000 & 1.000 & $\approx 6$--15 \\
50\% & 60/60 & 1.000 & 1.000 & $\approx 10$--25 \\
\bottomrule
\end{tabular}
\end{table}

\begin{table}[H]
\centering
\caption{F1 score vs.\ $z$-score threshold at 50\% mutation rate}
\label{tab:f1_zscore}
\begin{tabular}{cccc}
\toprule
$z$-score Threshold & Precision & Recall & F1 \\
\midrule
2 & 0.742 & 1.000 & 0.852 \\
3 & 0.891 & 1.000 & 0.942 \\
4 & 0.967 & 1.000 & 0.983 \\
5 & 0.992 & 1.000 & 0.996 \\
6 & 1.000 & 1.000 & 1.000 \\
7 & 1.000 & 1.000 & 1.000 \\
8 & 1.000 & 0.983 & 0.992 \\
\bottomrule
\end{tabular}
\end{table}

\noindent
\textbf{Planted motif recovery} (1\,Mb synthetic locus, 50\% background mutation):
canonical motif (position 500,000\,bp) recovered at rank 1 with $\rho = 0.9969$;
10\%-mutated copy (position 800,000\,bp) recovered at rank 2 with $\rho = 0.8979$.

%------------------------------------------------------------------------------
\section*{B.5\quad Cancer Suppression and Peto's Paradox}
\label{app:cancer_table}
\addcontentsline{toc}{section}{B.5\enspace Cancer Suppression}
%------------------------------------------------------------------------------

\noindent
Source: Chapter~\ref{chap:petos_paradox} and \citet{Sachikonye2025_Peto}.
The framework predicts cancer suppression factor $\Gamma_{\mathrm{error}} \propto \Omega^{-1/3}$
where $\Omega = N_{\mathrm{cells}} \times C_{\mathrm{val}}$ is the configuration
space volume. $\Gamma_{\mathrm{pred}}$ is normalized to human; $R_{\mathrm{obs}}$
is the observed lifetime cancer risk.

\vspace{6pt}

\begin{table}[H]
\centering
\caption{Cancer suppression: predicted vs.\ observed ($r = 0.94$)}
\label{tab:cancer_suppression}
\begin{tabular}{lcccc}
\toprule
Species & $N_{\mathrm{cells}}$ & $\Omega$ (arb.) &
$\Gamma_{\mathrm{pred}}$ & $R_{\mathrm{obs}}$ (lifetime, \%) \\
\midrule
\textit{Mus musculus}    & $3 \times 10^9$   & $7.8 \times 10^9$   & 0.22 & 40 \\
\textit{Rattus norvegicus} & $5\times10^9$  & $1.3\times10^{10}$  & 0.19 & 35 \\
\textit{Canis lupus}     & $3 \times 10^{11}$ & $7.8\times10^{11}$ & 0.07 & 27 \\
\textit{Homo sapiens}    & $3.7\times10^{13}$ & $1.2\times10^{14}$ & 0.02 & 39$^*$ \\
\textit{Balaenoptera musculus} & $10^{17}$ & $3.2\times10^{17}$  & 0.003 & $<1$ \\
\bottomrule
\multicolumn{5}{l}{$^*$Age-adjusted; unadjusted lifetime risk for long-lived humans.}
\end{tabular}
\end{table}

\noindent
Peto's paradox (large animals do not have proportionally higher cancer rates)
is reproduced by the $\Omega^{-1/3}$ suppression: the blue whale has
$\Gamma_{\mathrm{pred}} \approx 0.003$, consistent with the observed near-zero
cancer incidence in cetaceans \citep{peto1975cancers}.

%------------------------------------------------------------------------------
\section*{B.6\quad Immunological Response Comparison}
\label{app:immune_table}
\addcontentsline{toc}{section}{B.6\enspace Immunological Comparison}
%------------------------------------------------------------------------------

\noindent
Source: Chapter~\ref{chap:categorical_immunology} and \citet{Sachikonye2025_Immunity}.
The framework predicts that the immune response magnitude $R$ is proportional to
the categorical distance $\dC$ between the antigen's S-coordinate and the
self-antigen baseline. $R_{\mathrm{viral}}$ is measured as serum antibody titer
(log scale); $R_{\mathrm{toxin}}$ is measured as cytokine release index.

\vspace{6pt}

\begin{table}[H]
\centering
\caption{Immunological response: viral vs.\ toxin antigens}
\label{tab:immune_response}
\begin{tabular}{lcccc}
\toprule
Antigen Class & $\dC$ (S-coord.) & $R_{\mathrm{measured}}$ (AU) &
$R_{\mathrm{pred}}$ (AU) & $p$-value \\
\midrule
Influenza H3N2   & 0.71 & $89 \pm 12$ & $83$  & $0.31$ \\
SARS-CoV-2 spike & 0.68 & $94 \pm 9$  & $80$  & $0.08$ \\
HIV gp120        & 0.58 & $67 \pm 15$ & $68$  & $0.94$ \\
Tetanus toxin    & 0.33 & $44 \pm 7$  & $39$  & $0.37$ \\
LPS endotoxin    & 0.41 & $51 \pm 11$ & $48$  & $0.74$ \\
Diphtheria toxin & 0.29 & $38 \pm 6$  & $34$  & $0.41$ \\
\midrule
\multicolumn{2}{l}{Viral (mean)} & $83 \pm 14$ & --- & --- \\
\multicolumn{2}{l}{Toxin (mean)} & $44 \pm 9$  & --- & --- \\
\multicolumn{2}{l}{Viral vs.\ toxin} & \multicolumn{2}{c}{---} & $< 0.001$ \\
\bottomrule
\end{tabular}
\end{table}

\noindent
The significantly higher response to viral antigens ($p < 0.001$, two-sample
$t$-test) is predicted by the framework: viral proteins have higher $\dC$ from
self because they encode foreign partition coordinates, whereas toxins are
often structurally similar to host proteins (low $\dC$, reduced immune
recognition).

%------------------------------------------------------------------------------
\section*{B.7\quad Metabolic Hierarchy Information Flux}
\label{app:metabolic_table}
\addcontentsline{toc}{section}{B.7\enspace Metabolic Information Flux}
%------------------------------------------------------------------------------

\noindent
Source: Chapter~\ref{chap:metabolic_computer} and \citet{Sachikonye2025_Metabolic}.
Total metabolic information flux $I_{\mathrm{total}}$ is computed from the
mutual information at each level of the three-level hierarchy (glycolysis /
TCA / electron transport). Values in bits per metabolic cycle.

\vspace{6pt}

\begin{table}[H]
\centering
\caption{Metabolic hierarchy information flux under four conditions}
\label{tab:metabolic_flux}
\begin{tabular}{lcccc}
\toprule
Condition & $I_{\mathrm{glycolysis}}$ & $I_{\mathrm{TCA}}$ &
$I_{\mathrm{ETC}}$ & $I_{\mathrm{total}}$ \\
\midrule
Healthy (basal)       & 2.41 & 2.51 & 2.37 & 7.29 \\
Metformin-treated     & 2.41 & 2.02 & 2.37 & 6.80 \\
Insulin-resistant     & 2.41 & 1.87 & 2.02 & 6.30 \\
Cancer (Warburg)      & 2.41 & 0.94 & 0.87 & 4.22 \\
\bottomrule
\end{tabular}
\end{table}

\noindent
The Warburg effect (aerobic glycolysis in cancer cells) manifests as a
reduction in $I_{\mathrm{TCA}}$ and $I_{\mathrm{ETC}}$ to $\approx 37\%$ and
$\approx 37\%$ of healthy values, while $I_{\mathrm{glycolysis}}$ is maintained.
The total reduction from $7.29$ to $4.22$ bits/cycle ($\Delta I = 3.07$ bits)
is a partition-theoretic diagnostic for the Warburg metabolic phenotype.
Metformin acts by reducing $I_{\mathrm{TCA}}$ slightly (inhibiting complex I)
while maintaining higher-level homeostasis; insulin resistance reduces both
$I_{\mathrm{TCA}}$ and $I_{\mathrm{ETC}}$ proportionally.
